\documentclass[11pt,a4paper,twocolumn]{article}
\usepackage[utf8]{inputenc}
\usepackage{amsmath, amssymb}
\usepackage{geometry}
\geometry{a4paper, margin=1in}
\usepackage{enumitem}

\renewcommand{\familydefault}{\sfdefault} 

\title{\textbf{Latihan Soal Fungsi Aljabar, Komposisi, dan Invers}}
\author{}
\date{}

\begin{document}

\maketitle
\vspace{-1.5cm}

\begin{enumerate}[label=\arabic*., itemsep=8pt]
    \item Jika $f(x) = 2x + 3$ dan $g(x) = x - 5$, tentukan $(f + g)(x)$.
    \item Diberikan $p(x) = x^2 - 4$ dan $q(x) = x + 2$. Tentukan $(p - q)(x)$.
    \item Jika $f(x) = 3x$ dan $g(x) = x^2 + 1$, tentukan $(f \cdot g)(x)$.
    \item Diketahui $f(x) = x^2 - 9$ dan $g(x) = x - 3$. Tentukan $\left(\frac{f}{g}\right)(x)$ beserta syarat daerah asalnya.
    \item Jika $f(x) = 2x^2 + 3x - 5$ dan $g(x) = x^2 - 2x + 1$, tentukan $(f + g)(x)$.
    \item Diberikan $f(x) = 5x - 2$ dan $g(x) = 3x + 4$. Tentukan nilai dari $(f - g)(2)$.
    \item Diketahui $f(x) = \sqrt{x+2}$ dan $g(x) = \sqrt{x-2}$. Tentukan fungsi $(f \cdot g)(x)$.
    \item Jika $f(x) = x^3$ dan $g(x) = x^2$, tentukan $\left(\frac{f}{g}\right)(x)$ untuk $x \neq 0$.
    \item Diberikan $f(x) = 2x + 1$ dan $g(x) = 3x - 2$. Tentukan persamaan fungsi $(f \circ g)(x)$.
    \item Dengan $f(x) = x^2$ dan $g(x) = x + 4$, tentukan persamaan fungsi $(g \circ f)(x)$.
    \item Jika $f(x) = \frac{1}{x}$ dan $g(x) = 2x + 3$, tentukan fungsi $(f \circ g)(x)$.
    \item Diketahui $f(x) = x - 1$ dan $g(x) = x^2 + 2x + 1$. Tentukan fungsi $(f \circ g)(x)$.
    \item Diberikan $f(x) = 3x - 1$, $g(x) = 2x$, dan $h(x) = x^2$. Tentukan persamaan $(f \circ g \circ h)(x)$.
    \item Jika $f(x) = 2x - 5$, tentukan nilai dari $(f \circ f)(3)$.
    \item Diketahui $f(x) = \frac{x+1}{x-2}$ dan $g(x) = 3x$. Tentukan nilai dari $(f \circ g)(1)$.
    \item Diketahui $(f \circ g)(x) = 2x + 5$ dan $f(x) = 2x + 1$. Tentukan persamaan fungsi $g(x)$.
    \item Jika $(g \circ f)(x) = x^2 - 4x + 5$ dan $f(x) = x - 2$, tentukan persamaan fungsi $g(x)$.
    \item Diberikan $(f \circ g)(x) = 6x - 3$ dan $g(x) = 2x - 1$. Tentukan persamaan fungsi $f(x)$.
    \item Diketahui $(f \circ g)(x) = \frac{2x+3}{x-1}$ dan $f(x) = x + 2$. Tentukan fungsi $g(x)$.
    \item Jika $(f \circ g)(x) = 4x^2 + 8x - 3$ dan $g(x) = 2x + 2$, tentukan persamaan fungsi $f(x)$.
    \item Diberikan $(g \circ f)(x) = 3x + 2$ dan $g(x) = 3x - 4$. Tentukan persamaan fungsi $f(x)$.
    \item Diketahui $(f \circ g)(x) = x^2 + 4x + 5$ dan $f(x) = x^2 + 1$. Tentukan $g(x)$ dengan asumsi $g(x)$ adalah fungsi linear bergradien positif.
    \item Tentukan rumus fungsi invers dari $f(x) = 3x - 7$.
    \item Jika $g(x) = \frac{2x + 1}{x - 3}$, tentukan fungsi invers $g^{-1}(x)$.
    \item Diberikan $h(x) = 5 - 2x$. Tentukan nilai dari $h^{-1}(1)$.
    \item Tentukan fungsi invers dari $f(x) = \frac{x}{x+2}$.
    \item Jika $f(x) = x^2 - 4$ untuk $x \geq 0$, tentukan fungsi invers $f^{-1}(x)$.
    \item Diberikan $f(x) = \sqrt{2x - 5}$. Tentukan fungsi invers $f^{-1}(x)$.
    \item Diketahui $f(x) = 2x + 1$ dan $g(x) = \frac{x-1}{3}$. Tentukan fungsi invers $(f \circ g)^{-1}(x)$.
    \item Jika $f(x) = 3x + 2$ dan $g(x) = 4x - 1$, tentukan $(g^{-1} \circ f^{-1})(x)$.
\end{enumerate}

\end{document}
